\begin{tabularx}{\textwidth}{L{34mm}X X}
\toprule
\textbf{Fogalom} & \textbf{Jelentése} & \textbf{Java-beli eszköz / megjegyzés} \\
\midrule
Öröklődés & Új osztály létrehozása meglévő osztály specializálásával. & \code{extends}; Java-ban osztályok között egyszeres öröklődés van, de több interfész is megvalósítható. \\
Ősosztály & Az általánosabb osztály, amely közös állapotot és viselkedést adhat. & Közös mezők, metódusok, konstruktorlogika; minden osztály közvetve a \code{Object} leszármazottja. \\
Leszármazott osztály & Az ősosztály speciálisabb változata. & Örökli az elérhető tagokat, új tagokat vehet fel, metódusokat felüldefiniálhat. \\
Polimorfizmus & Ugyanazon típusú referencia mögött többféle tényleges objektum állhat. & Őstípusú, absztrakt osztály típusú vagy interfész típusú referencia; dinamikus metóduskötés. \\
Túlterhelés & Azonos nevű metódusok eltérő paraméterlistával. & Fordítási időben dől el, hogy melyik metódusváltozat illeszkedik. \\
Felüldefiniálás & Örökölt példánymetódus új implementációja leszármazottban. & \code{@Override}; futási időben a dinamikus típus alapján választódik ki. \\
\bottomrule
\end{tabularx}
